%%%%%解析几何模版%%%%%
\chapter*{解析几何每日一题}

\refer{《解析几何》}

\exercise{一平面通过两点$M_1\qty(1,1,1),\ M_2\qty(0,1,-1)$且垂直于平面$x+y+z=0$，求它的方程.}

\teacher{马心怡}

\postdate{2023年4月19日}

\result{2x-y-z=0}

\analysis{本题考查的是向量运算以及对于法向量的应用，在面对求平面以及直线方 程的时候，我们应该将其与法向量以及方向向量联系起来，通过法向量与 方向向量再加上额外的点可以生成唯一平面，而对于求法向量可以利用点 乘和叉乘快速求解.

提示：平面方程$Ax+By+Cz=0$的法向量为$\qty(A,B,C)$}

\solution{
    设平面的法向量为$\vv{n}=\qty(A,B,C)$且$\vv{M_1M_2}=\qty(-1,0,2)$为所求向量，则
    \formula{\vv{n}\cdot\vv{M_1M_2}=-A-2C=0}
    又因为所求平面垂直于已知平面
    \formula{x+y+z=0}
    所以两平面的法向量也相互垂直，即
    \formula{A+B+C=0}
    联立得
    \formula{A=-2C,\ B=C}
    则有
    \formula{A\qty(x-1)+B\qty(y-1)+C\qty(z-1)=0}
    代入求得
    \formula{22x-y-z=0}
}
